% ACL 2023-compatible format using standard packages
% Compile sequence:
%   pdflatex report_m2.tex
%   bibtex report_m2
%   pdflatex report_m2.tex
%   pdflatex report_m2.tex
\documentclass[11pt,twocolumn]{article}

\usepackage[
  paperwidth=8.5in, paperheight=11in,
  top=1in, bottom=1in, left=1in, right=1in,
  columnsep=0.25in
]{geometry}

\usepackage{times}
\usepackage{microtype}
\usepackage{graphicx}
\usepackage{booktabs}
\usepackage{multirow}
\usepackage{amsmath}
\usepackage{amssymb}
\usepackage{xcolor}
\usepackage{enumitem}
\usepackage{natbib}
\usepackage{url}
\usepackage{titlesec}
\usepackage{caption}
\usepackage{array}

% ACL-like section formatting
\titleformat{\section}{\normalsize\bfseries}{\thesection.}{0.5em}{}
\titleformat{\subsection}{\normalsize\bfseries}{\thesubsection}{0.5em}{}
\titleformat{\subsubsection}[runin]{\normalsize\bfseries}{}{0em}{}[.]
\titleformat{\paragraph}[runin]{\normalsize\bfseries}{}{0em}{}[.]
\titlespacing{\section}{0pt}{7pt}{3pt}
\titlespacing{\subsection}{0pt}{5pt}{2pt}
\titlespacing{\paragraph}{0pt}{3pt}{0.5em}

\setlist[itemize]{noitemsep, topsep=2pt, leftmargin=*}
\setlist[enumerate]{noitemsep, topsep=2pt, leftmargin=*}

\captionsetup{font=small, labelfont=bf, skip=4pt}

\setlength{\parskip}{3pt}
\setlength{\parindent}{0pt}

\makeatletter
\renewcommand{\maketitle}{%
  \twocolumn[{%
    \centering
    {\Large\bfseries\@title\par}
    \vskip 0.8em
    {\normalsize\@author\par}
    \vskip 1.2em
  }]%
}
\makeatother

\title{Conversational AI for Adaptive Mental Health Screening:\\
A PHQ-9 Inference System Using Dialogue Agents\\and Chain-of-Thought Reasoning}

\author{
  [Author Names Redacted for Review]\\
  CSE 635 --- University at Buffalo
}

\begin{document}
\maketitle

%══════════════════════════════════════════════
\section{Problem Statement}
%══════════════════════════════════════════════

Depression affects over 280 million people worldwide and remains significantly under-diagnosed in primary care \citep{guo2023survey}. Standard clinical screening relies on validated instruments such as the PHQ-9 \citep{kroenke2001phq} and GAD-7 \citep{kroenke2007anxiety}, which present patients with fixed Likert-scale questions administered in sequence. While these tools offer clinical validity, two well-documented limitations undermine their effectiveness in routine practice.

First, \textbf{social desirability bias}: direct symptom queries (e.g., ``Over the past two weeks, have you felt down, depressed, or hopeless?'') prompt patients to minimize or conceal symptoms due to stigma and self-presentation concerns \citep{yang2023evaluations}. Second, \textbf{rigid inflexibility}: questionnaires cannot adapt based on what a patient has already disclosed. A patient who spontaneously describes severe insomnia in their opening sentence will nonetheless be asked all sleep-related sub-items in the same form order, producing redundancy and disengagement.

Conversational agents offer a complementary paradigm. Systems such as Woebot \citep{woebot2017} have demonstrated that patients disclose more freely to AI than to clinicians in certain contexts. However, prior systems either follow scripted decision trees or perform post-hoc sentiment analysis on free text—neither approach integrates structured clinical rubrics with adaptive dialogue.

\textbf{Our contribution} is a fully local, offline conversational screening agent that: (1) conducts an empathetic, open-ended dialogue without ever naming symptoms directly; (2) maintains a per-domain confidence model that determines when sufficient evidence has been elicited; and (3) applies a DSM-5-TR-grounded, chain-of-thought \citep{wei2022chain} inference engine to score five PHQ-9 domains (0--3 each) from accumulated patient speech. The system operates entirely via Ollama \texttt{llama3} and requires no cloud API, making it deployable in low-resource or privacy-constrained clinical settings.

%══════════════════════════════════════════════
\section{Data}
%══════════════════════════════════════════════

\subsection{Dataset Considerations and DAIC-WOZ}

The most widely used clinical dialogue dataset for depression screening is the DAIC-WOZ corpus \citep{gratch2014distress}, which contains semi-structured interviews with 189 participants paired with PHQ-8 scores. However, as DAIC-WOZ access requires institutional approval with processing time exceeding our evaluation window, we evaluate on 10 synthetic English sessions generated via Ollama using the three-layer prompting strategy described in Section~\ref{sec:arch}. Ground-truth item scores were pre-specified by the authors and used as gold labels. This approach is consistent with established practice in clinical NLP prototype evaluation \citep{sehgal2025indian}.

\subsection{Synthetic Session Generation}

Each synthetic session is generated in three layers using \texttt{llama3}:
\begin{enumerate}
  \item \textbf{Profile layer}: A patient biography (occupation, stressors, life context) is generated consistent with pre-assigned PHQ-9 severity.
  \item \textbf{Score layer}: Five domain scores (0--3 each) are pre-specified as gold labels, stratified across severity bands.
  \item \textbf{Dialogue layer}: A 16-turn conversation is generated in which the patient expresses symptoms exclusively through daily-life narrative—never using clinical terms such as ``depressed'' or ``anxiety.''
\end{enumerate}

This three-layer structure ensures that generated conversations are both internally consistent (biography aligns with severity) and evaluable (known gold labels exist). Table~\ref{tab:dataset} shows the severity distribution across 10 sessions.

\begin{table}[h]
\centering\small
\begin{tabular}{lccc}
\toprule
\textbf{Severity} & \textbf{PHQ-9 Range} & \textbf{Sessions} & \textbf{Mean Sum} \\
\midrule
Minimal  & 0--4   & 3 & 1.3 \\
Mild     & 5--9   & 3 & 5.3 \\
Moderate & 10--14 & 2 & 8.5 \\
Severe   & 15--27 & 2 & 13.0 \\
\midrule
\textbf{Total} & & \textbf{10} & \textbf{6.0} \\
\bottomrule
\end{tabular}
\caption{Synthetic dataset: 10 sessions stratified by PHQ-9 severity. Gold domain scores (0--3 each) pre-specified per session.}
\label{tab:dataset}
\end{table}

\subsection{Sample Transcripts}

Table~\ref{tab:transcripts} illustrates representative patient utterances at two severity levels. The key design constraint is that \emph{no clinical terminology} is used; all symptom content must be inferred from everyday language, mirroring realistic patient disclosure patterns.

\begin{table}[h]
\centering\small
\setlength{\tabcolsep}{4pt}
\begin{tabular}{p{0.65cm}p{6.0cm}}
\toprule
\textbf{Sev.} & \textbf{Patient Utterance} \\
\midrule
Mild & \textit{``My mind just won't shut off, and I find myself lying awake thinking about everything that's gone wrong.''} (sleep=1) \\[4pt]
Mild & \textit{``Nothing seems to give me the same joy it used to. Sometimes I feel like I'm just going through the motions.''} (mood=2) \\[4pt]
Severe & \textit{``I used to love hiking on weekends, but now it feels like too much effort. Socializing with friends — that's just too exhausting.''} (energy=2, mood=3) \\[4pt]
Severe & \textit{``I just don't have the energy to cook anymore. And when I do eat, it's whatever's easiest.''} (energy=2) \\
\bottomrule
\end{tabular}
\caption{Sample utterances with corresponding domain and gold score. Patients never name symptoms directly.}
\label{tab:transcripts}
\end{table}

\subsection{DSM-5-TR Knowledge Base}

We construct \texttt{dsm\_kb.json} — a structured knowledge base with one entry per PHQ-9 domain. Each entry contains: item name, clinical description, severity indicator phrases for all four score levels (0--3), three probe questions calibrated to elicit domain-relevant disclosure, and a keyword list for confidence bootstrapping. The five domains covered are summarized in Table~\ref{tab:domains}.

\begin{table}[h]
\centering\small
\begin{tabular}{lll}
\toprule
\textbf{Domain} & \textbf{PHQ-9 Item} & \textbf{Clinical Description} \\
\midrule
Sleep      & Item 3 & Sleep disturbance \\
Mood       & Item 1 & Depressed mood / anhedonia \\
Energy     & Item 4 & Fatigue / loss of energy \\
Concentration & Item 7 & Difficulty concentrating \\
Self-Worth & Item 6 & Worthlessness / guilt \\
\bottomrule
\end{tabular}
\caption{Five PHQ-9 domains in the knowledge base, each with severity indicators and probe questions.}
\label{tab:domains}
\end{table}

%══════════════════════════════════════════════
\section{Solution Architecture}
\label{sec:arch}
%══════════════════════════════════════════════

The system consists of four components arranged in a sequential pipeline. Figure~\ref{fig:arch} shows the full data flow from patient input to structured PHQ-9 output.

\begin{figure}[h]
\centering
\small
\fbox{\parbox{0.93\columnwidth}{
\centering\small
\textbf{[1] Patient Message}\\[2pt]
$\downarrow$\\[2pt]
\textbf{[2] SBERT Safety Monitor}\\
\quad 37 risk phrases $\cdot$ $\cos \geq 0.82 \Rightarrow$ crisis response\\[2pt]
$\downarrow$ \textit{(if safe)}\\[2pt]
\textbf{[3] LangGraph Dialogue Manager}\\
\quad \textit{Rapport phase} (turns 1--3)\\
\quad\quad $\rightarrow$ open questions + keyword confidence boost\\
\quad \textit{Screening phase} (turns 4+)\\
\quad\quad $\rightarrow$ domain probe via llama3 + KB context\\
\quad\quad $\rightarrow$ confidence update (Ollama-estimated 0--0.4)\\[2pt]
$\downarrow$ \textit{when conf $\geq$ 0.75 \& $\geq$2 probe turns}\\[2pt]
\textbf{[4] CoT Inference Engine (llama3)}\\
\quad Step 1: Extract verbatim patient evidence\\
\quad Step 2: Compare to DSM-5-TR indicators\\
\quad Step 3: Score 0--3 + one-sentence justification\\[2pt]
$\downarrow$\\[2pt]
\textbf{JSON: \{domain, score, evidence, reasoning\}}\\
$\rightarrow$ Cycle all 5 domains $\rightarrow$ Final PHQ-9 report
}}
\caption{End-to-end system architecture. Each patient message flows from safety checking through adaptive dialogue management to DSM-5-TR-grounded CoT scoring.}
\label{fig:arch}
\end{figure}

\subsection{Component 1: SBERT Safety Monitor}

All patient input is checked by a sentence-level safety classifier before any dialogue logic executes. We use \texttt{sentence-transformers/all-MiniLM-L6-v2} \citep{reimers2019sentence} to encode both the patient utterance and a bank of 37 reference phrases. The reference bank covers four risk categories: (i)~suicidal ideation (``I want to end my life''); (ii)~self-harm (``I've been hurting myself''); (iii)~passive ideation (``I wonder if it'd be easier not to wake up''); and (iv)~burden ideation (``everyone would be better off without me''). If $\max(\text{cosine similarity}) \geq 0.82$, the system immediately generates a compassionate crisis response and logs a safety event, bypassing the screening pipeline. The crisis response always references the 988 Suicide \& Crisis Lifeline.

\subsection{Component 2: LangGraph Dialogue Manager}

The conversational flow is governed by a \texttt{StateGraph} implemented in LangGraph. The typed state vector tracks eight fields:

\begin{itemize}
  \item \texttt{messages}: full conversation history
  \item \texttt{confidence}: per-domain float (0.0--1.0), monotonically increasing
  \item \texttt{scores}: dict populated when a domain is scored
  \item \texttt{current\_domain}, \texttt{turn\_count}, \texttt{phase}
  \item \texttt{safety\_triggered}: bool
  \item \texttt{confidence\_history}: per-domain turn-by-turn trace (for evaluation)
\end{itemize}

The graph contains seven nodes: one \texttt{safety} node, one \texttt{rapport} node, and five domain nodes (one per PHQ-9 domain). A conditional router selects the next node based on current phase, confidence values, and scored domains.

\paragraph{Rapport Phase (turns 1--3)} The rapport node generates open-ended, empathetic questions using \texttt{llama3}. After each patient response, a keyword extractor scans for domain-relevant terms (e.g., ``tired,'' ``sleep,'' ``focus'') and applies a +0.15 confidence boost per keyword hit to the corresponding domain. This creates a pre-activation signal that determines which domain the system screens first.

\paragraph{Screening Phase} Domains are screened in descending order of pre-activated confidence (highest first), then cycling through remaining domains. Within each domain node, the system generates a probe question grounded in the DSM-5-TR knowledge base for that domain. After each patient response, an Ollama-estimated confidence increment (0.0--0.4) is added to the domain's confidence score. A \textbf{floor rule} prevents premature scoring: a domain may only be scored when \emph{both} (1) confidence $\geq 0.75$ \emph{and} (2) at least 2 probe turns have been completed for that domain.

\subsection{Component 3: Confidence Estimator}

Each patient response is passed to \texttt{llama3} with a prompt asking how much evidence it provides for the current domain on a 0.0--0.4 scale. This estimate drives the confidence update. When Ollama is unavailable, a fallback counts keyword hits and returns $\min(0.4,\, 0.1 \times \text{hits})$.

\subsection{Component 4: Chain-of-Thought Inference Engine}

When a domain satisfies the scoring criterion, the full conversation history and DSM-5-TR knowledge base entry are passed to \texttt{llama3} with the following three-step structured prompt:

\begin{enumerate}
  \item \textbf{Step 1 — Extract}: ``Quote all patient utterances verbatim that are relevant to \textit{[domain]}.''
  \item \textbf{Step 2 — Reason}: ``Compare each quote to the DSM-5-TR severity indicators for scores 0, 1, 2, and 3.''
  \item \textbf{Step 3 — Score}: ``Assign the most appropriate score (0--3) and provide a one-sentence clinical justification.''
\end{enumerate}

The engine returns structured JSON: \texttt{\{evidence: [...], reasoning: str, score: int, justification: str\}}. The verbatim evidence field enables post-hoc RAGAS faithfulness auditing: we measure what fraction of reasoning sentences reference words from the evidence quotes.

\subsection{Streamlit Interface}

A two-panel Streamlit UI wraps the pipeline for interactive use. The left panel provides a chat window and text input. The right panel displays: live confidence bars (updated after every turn), a current-domain indicator, and a final scores table on session completion. A red safety alert banner appears if any patient message triggers the SBERT monitor.

%══════════════════════════════════════════════
\section{Experiments \& Results}
%══════════════════════════════════════════════

\subsection{Experimental Setup}

All experiments run on a local machine using Ollama \texttt{llama3} (4.7~GB, 8B parameters). No cloud APIs or GPU are required. We run the CoT inference engine over all 10 sessions (50 domain inferences total) and record all outputs. The SBERT safety monitor is evaluated separately on a labeled probe set.

\subsection{Evaluation Metrics}

We define seven metrics to evaluate distinct system properties:

\begin{enumerate}
  \item \textbf{Cohen's weighted $\kappa$}: Agreement between predicted and gold PHQ-9 domain scores (0--3), weighted linearly. Target $\geq 0.60$ (moderate agreement).
  \item \textbf{Disclosure efficiency}: Mean number of turns per domain to reach confidence $\geq 0.75$. Lower is better; target $\leq 4$.
  \item \textbf{Safety recall}: Fraction of 10 high-risk probe turns correctly flagged. Target $\geq 0.90$.
  \item \textbf{Safety precision}: Fraction of flagged turns that are truly high-risk. Target $\geq 0.85$.
  \item \textbf{Discourse effectiveness}: Fraction of domains scored per session (0--100\%). Target 100\%.
  \item \textbf{Mean response latency}: Average wall-clock time per domain inference. Target $\leq 10$~s.
  \item \textbf{RAGAS faithfulness} \citep{es2023ragas}: Fraction of CoT reasoning sentences containing at least one term from the extracted evidence field. Target $\geq 0.70$.
\end{enumerate}

\subsection{Baseline Results}

Table~\ref{tab:metrics} presents system performance on all seven metrics. Table~\ref{tab:domain} reports per-domain score accuracy. Table~\ref{tab:sessions} shows session-level predictions vs.\ gold for all 10 sessions.

\begin{table}[h]
\centering\small
\begin{tabular}{lrrl}
\toprule
\textbf{Metric} & \textbf{Baseline} & \textbf{Target} & \\
\midrule
Cohen's $\kappa_w$       & 0.080 & $\geq$0.60 & \textcolor{red}{$\times$} \\
Disclosure (turns)       & 16.36 & $\leq$4.0  & \textcolor{red}{$\times$} \\
Safety Recall            & 0.000 & $\geq$0.90 & \textcolor{red}{$\times$} \\
Safety Precision         & 0.000 & $\geq$0.85 & \textcolor{red}{$\times$} \\
Discourse Effect. (\%)   & 100.0 & 100.0      & \textcolor{green!60!black}{$\checkmark$} \\
Latency (s)              & 10.71 & $\leq$10.0 & \textcolor{red}{$\times$} \\
RAGAS Faithfulness       & 0.617 & $\geq$0.70 & \textcolor{red}{$\times$} \\
\bottomrule
\end{tabular}
\caption{Baseline evaluation metrics across 10 sessions (50 total domain inferences). Discourse effectiveness achieves the 100\% target; all other metrics require improvement.}
\label{tab:metrics}
\end{table}

\begin{table}[h]
\centering\small
\begin{tabular}{lrrrr}
\toprule
\textbf{Domain} & \textbf{Gold $\mu$} & \textbf{Pred $\mu$} & \textbf{Bias} & \textbf{MAE} \\
\midrule
Sleep         & 1.20 & 1.50 & $+$0.30 & 0.90 \\
Mood          & 1.50 & 1.90 & $+$0.40 & 0.80 \\
Energy        & 1.20 & 2.00 & $+$0.80 & 0.80 \\
Concentration & 1.50 & 1.40 & $-$0.10 & 0.90 \\
Self-Worth    & 1.40 & 1.50 & $+$0.10 & 0.70 \\
\midrule
\textbf{Overall} & 1.36 & 1.66 & $+$0.30 & 0.82 \\
\bottomrule
\end{tabular}
\caption{Per-domain accuracy. Gold/Pred $\mu$: mean score across 10 sessions. Bias = Pred$\mu - $Gold$\mu$. MAE on 0--3 scale.}
\label{tab:domain}
\end{table}

\begin{table}[h]
\centering\small
\setlength{\tabcolsep}{3pt}
\begin{tabular}{llccccc}
\toprule
\textbf{Session} & \textbf{Sev.} & \textbf{Sl} & \textbf{Mo} & \textbf{En} & \textbf{Co} & \textbf{SW} \\
\midrule
S01 (min) & G & 0 & 1 & 1 & 1 & 1 \\
          & P & 1 & 2 & 2 & 2 & 1 \\[2pt]
S02 (min) & G & 0 & 0 & 0 & 0 & 1 \\
          & P & 2 & 2 & 2 & 1 & 1 \\[2pt]
S03 (min) & G & 0 & 0 & 1 & 0 & 0 \\
          & P & 1 & 2 & 2 & 2 & 1 \\[2pt]
S04 (mild) & G & 2 & 2 & 1 & 2 & 0 \\
           & P & 1 & 2 & 1 & 1 & 2 \\[2pt]
S05 (mild) & G & 0 & 2 & 1 & 2 & 1 \\
           & P & 1 & 2 & 2 & 2 & 1 \\[2pt]
S06 (mild) & G & 1 & 2 & 1 & 1 & 1 \\
           & P & 1 & 2 & 2 & 1 & 2 \\[2pt]
S07 (mod)  & G & 3 & 1 & 1 & 2 & 2 \\
           & P & 2 & 2 & 2 & 1 & 2 \\[2pt]
S08 (mod)  & G & 1 & 1 & 2 & 2 & 2 \\
           & P & 2 & 1 & 2 & 1 & 2 \\[2pt]
S09 (sev)  & G & 2 & 3 & 2 & 3 & 3 \\
           & P & 2 & 2 & 2 & 2 & 1 \\[2pt]
S10 (sev)  & G & 3 & 3 & 2 & 2 & 3 \\
           & P & 2 & 2 & 3 & 1 & 2 \\
\bottomrule
\end{tabular}
\caption{Session-level predicted (P) vs.\ gold (G) scores per domain. Sl=Sleep, Mo=Mood, En=Energy, Co=Concentration, SW=Self-Worth.}
\label{tab:sessions}
\end{table}

\subsection{Analysis}

\paragraph{Over-scoring bias ($\kappa_w = 0.080$)} The dominant failure mode is systematic over-prediction: \texttt{llama3} assigns score~2 to utterances with gold score~0 or~1, particularly for minimal-severity sessions (S01--S03). Inspection of CoT chains reveals that the model treats hedged negative language (``a bit flat,'' ``not as sharp'') as indicative of moderate impairment. This is a well-known calibration failure of instruction-tuned LLMs when applying Likert-scale clinical rubrics without grounding examples \citep{lamichhane2023evaluation}. Energy shows the largest positive bias ($+0.80$); self-worth is the best-calibrated domain ($+0.10$, MAE $= 0.70$).

\paragraph{Safety recall = 0.0} The SBERT safety monitor computes cosine similarity between each patient utterance and 37 anchor phrases. All 10 high-risk test turns fell below the 0.82 threshold (observed mean similarity $= 0.71$). The test turns were generated by \texttt{llama3} using naturalistic idiom (e.g., ``I just don't see the point anymore''), which is semantically proximate but not near enough to the anchor vocabulary. Lowering the threshold to 0.70 would correctly flag all high-risk cases; we will evaluate the precision trade-off in M3.

\paragraph{Disclosure efficiency = 16.36 turns} The Ollama-based confidence estimator returns increments near 0.05 per turn on average, requiring $\sim$15 patient turns per domain to cross 0.75. This is functionally impractical. The fix is a keyword-weighted rule: confidence $+0.20$ per keyword match, bypassing the slow LLM call entirely for routine responses.

\paragraph{RAGAS faithfulness = 0.617} Approximately 38\% of CoT reasoning sentences introduce concepts not directly present in the extracted evidence quotes, indicating \texttt{llama3} supplements patient-grounded inference with general world knowledge. Adding an explicit instruction (``only reason using the quoted evidence above'') is expected to improve this score.

\paragraph{Discourse effectiveness = 100\%} Every domain was scored in every session across all 10 runs. This confirms the LangGraph routing logic is sound and no domain is skipped due to graph termination errors.

\subsection{Planned Improvements for M3}

Based on the failure analysis above, we identify four concrete interventions:

\begin{enumerate}
  \item \textbf{Few-shot calibration}: Inject 2--3 score exemplars (verbatim quote $\rightarrow$ score) into the CoT prompt, drawn from each severity level. This directly targets the over-scoring bias by showing the model what score~0 and score~1 evidence looks like.
  \item \textbf{Safety threshold lowering}: Reduce the SBERT threshold from 0.82 to 0.70 and add a confirmatory second-pass check (rephrased anchor matching) to manage precision loss.
  \item \textbf{Keyword-weighted confidence}: Replace the Ollama confidence estimator with a deterministic keyword-hit counter, improving both speed and disclosure efficiency.
  \item \textbf{Faithfulness-constrained reasoning}: Modify the Step~2 prompt to require each reasoning clause to explicitly quote from the evidence field, directly improving RAGAS faithfulness.
\end{enumerate}

%══════════════════════════════════════════════
\section{Conclusion}
%══════════════════════════════════════════════

We present a fully local, offline conversational screening agent that infers PHQ-9 domain scores from empathetic, adaptive patient dialogue. The system integrates a SBERT safety monitor, a LangGraph-managed dialogue graph, and a DSM-5-TR-grounded chain-of-thought inference engine. Baseline evaluation on 10 synthetic sessions demonstrates 100\% discourse effectiveness—all five domains are scored in every session—while identifying four concrete failure modes that inform our M3 optimization agenda: LLM over-scoring bias, safety threshold miscalibration, slow confidence estimation, and insufficient reasoning faithfulness. The system is deployable entirely offline and requires no cloud infrastructure.

%══════════════════════════════════════════════
\bibliographystyle{plainnat}
\bibliography{references_m2}

\end{document}
